\section{Optional garbage-collection layer}
\label{sec:garbage-collection}

This section is an optional conservative extension of the VM semantics.  It is
kept in this separate source file, and no preceding rule depends on one of its
auxiliary definitions.  Omitting the single
\verb|\section{Optional garbage-collection layer}
\label{sec:garbage-collection}

This section is an optional conservative extension of the VM semantics.  It is
kept in this separate source file, and no preceding rule depends on one of its
auxiliary definitions.  Omitting the single
\verb|\section{Optional garbage-collection layer}
\label{sec:garbage-collection}

This section is an optional conservative extension of the VM semantics.  It is
kept in this separate source file, and no preceding rule depends on one of its
auxiliary definitions.  Omitting the single
\verb|\section{Optional garbage-collection layer}
\label{sec:garbage-collection}

This section is an optional conservative extension of the VM semantics.  It is
kept in this separate source file, and no preceding rule depends on one of its
auxiliary definitions.  Omitting the single
\verb|\input{gc-semantics}| line restores the monotonically growing heaps of
Sections~3--11.  Including it replaces that implementation abstraction with an
explicit, nondeterministic collector.  The collector introduces no Newspeak
syntax and never runs inside an actor step or reflective transaction.

The construction follows the explicit-heap rewrite method of Morrisett,
Felleisen, and Harper, \emph{Abstract Models of Memory Management} (1995), and
the treatment of observable collection in Donnelly, Hallett, and Kfoury,
\emph{Formal Semantics of Weak References} (2006).  The distinctive Newspeak
observations are Equation~\eqref{eq:all-instances-of}, whose implicit instance
registry is weak, and the platform classes \kw{WeakArray} and \kw{WeakMap}.
Collection can therefore change both later population snapshots and later
reads of weak containers.

\subsection{The collectable store}

Garbage collection applies to ordinary heap records and to the routing records
of promise and far-reference handles.  For
$\Omega=\langle P,V,\mathcal H,\mathcal X,\mathcal Q,\mathcal N,
\Pi,\Phi,\Gamma,\mathcal E,\mathcal D,w,n\rangle$, define the disjoint reference
store
\begin{equation}\label{eq:gc-reference-store}
 \mathsf{refStore}(\Omega)
   =\mathsf{vmStore}(\Omega)\uplus\Pi\uplus\Phi.
\end{equation}
The union is disjoint because its three domains are respectively subsets of
$\mathit{ObjRef}$, $\mathit{PromiseId}$, and $\mathit{FarRefId}$.  Put
$\mathcal S_\Omega=\mathsf{refStore}(\Omega)$ and expose a uniform payload
projection:
\begin{equation}\label{eq:gc-reference-payload}
 \mathsf{payload}_{\Omega}(x)=
 \begin{cases}
  \mathsf{vmStore}(\Omega)(x) & x\in\mathit{ObjRef},\\
  \Pi(x) & x\in\mathit{PromiseId},\\
  \Phi(x) & x\in\mathit{FarRefId}.
 \end{cases}
 \qquad x\in\dom(\mathcal S_\Omega).
\end{equation}
Thus a live promise strongly retains its settlement value and registered
waiters, and a live far-reference handle strongly retains its near target.
The tables $\Pi$ and $\Phi$ do not retain their own keys merely by containing
them; otherwise every handle and every remote target would live forever.

\subsection{WeakArray and WeakMap}
\label{sec:weak-containers}

\kw{WeakArray} and \kw{WeakMap} are mutable platform objects and therefore
belong to exactly one actor heap.  Their concrete VM layouts are abstracted by
two additional object-record forms:
\begin{equation}\label{eq:weak-container-records}
\begin{aligned}
 H_A(a_w)&=\kw{weakArray}\langle C,\langle y_1,\ldots,y_m\rangle\rangle,
   &a_w&\in\mathit{WeakArrayId},\\
 H_A(m_w)&=\kw{weakMap}\langle C,E_w\rangle,
   &m_w&\in\mathit{WeakMapId},
\end{aligned}
\end{equation}
where $\mathit{WeakArrayId}$ and $\mathit{WeakMapId}$ are disjoint subsets of
$\mathit{ObjRef}$, every $y_i\in\mathit{Ref}$, and
$E_w:\mathit{ObjRef}\rightharpoonup\mathit{Ref}$ is finite.  $C$ and any
ordinary fixed fields of the implementation are strong; the sequence elements
and map keys are weak.  A map value is conditionally strong exactly while its
key is strongly reachable.  This is the usual ephemeron interpretation of a
weak-key map: a value-to-key cycle cannot keep the key alive.

Every stored weak identity initially resolves; weakness controls retention, not
whether an arbitrary dangling identity may be written:
\begin{equation}\label{eq:weak-container-validity}
\begin{aligned}
 \mathsf{weakArrayStoreOK}(\Omega)
 &\Longleftrightarrow
   \forall A,a_w,C,\bar y.\
     \mathcal H(A)(a_w)=\kw{weakArray}\langle C,\bar y\rangle\\[-1mm]
 &\hspace{35mm}\Longrightarrow
     \mathsf{refs}(\bar y)\subseteq\dom(\mathcal S_\Omega),\\
 \mathsf{weakMapStoreOK}(\Omega)
 &\Longleftrightarrow
   \forall A,m_w,C,E_w.\
     \mathcal H(A)(m_w)=\kw{weakMap}\langle C,E_w\rangle\\[-1mm]
 &\hspace{35mm}\Longrightarrow
     \dom(E_w)\cup\mathsf{refs}(\mathsf{range}(E_w))
       \subseteq\dom(\mathcal S_\Omega),\\
 \mathsf{weakContainersOK}(\Omega)
 &\Longleftrightarrow
   \mathsf{weakArrayStoreOK}(\Omega)\land\mathsf{weakMapStoreOK}(\Omega).
\end{aligned}
\end{equation}

The platform method contract is given by distinct partial store functions
rather than overloaded notation:
\begin{equation}\label{eq:weak-array-operations}
\begin{aligned}
 \mathsf{weakArrayRead}(H_A,a_w,i)&=y_i,\\
 \mathsf{writeWeakArrayCell}(H_A,a_w,i,v)
   &=H_A[a_w\mapsto
       \kw{weakArray}\langle C,\bar y[i\mapsto v]\rangle],\\
 &\hspace{35mm}1\le i\le |\bar y|.
\end{aligned}
\end{equation}
\begin{equation}\label{eq:weak-map-operations}
\begin{aligned}
 \mathsf{weakMapLookup}(H_A,m_w,k)&=
   \begin{cases}E_w(k)&k\in\dom(E_w),\\ \nil&k\notin\dom(E_w),\end{cases}\\
 \mathsf{writeWeakMapEntry}(H_A,m_w,k,v)
   &=H_A[m_w\mapsto\kw{weakMap}\langle C,E_w[k\mapsto v]\rangle].
\end{aligned}
\end{equation}
Here $H_A(a_w)$ and $H_A(m_w)$ have the respective forms of
Equation~\eqref{eq:weak-container-records}.  Allocation creates an all-$\nil$
array of the requested nonnegative size or a map with $E_w=\varnothing$, using
the ordinary global-freshness obligation of Equation~\eqref{eq:gc-wrapper-state}.
The two writes are defined only when their new key and value references resolve
in $\mathcal S_\Omega$, and they preserve Equation~
\eqref{eq:weak-container-validity}.
Out-of-bounds array indices are undefined and are reified by the platform
exception boundary.  Storing $\nil$ is permitted.  Since lookup also returns
$\nil$ for an absent key, the public \kw{WeakMap} protocol deliberately does
not distinguish absence from a stored $\nil$ value.

The state outside the collectable store is
\begin{equation}\label{eq:gc-control-state}
 \mathsf{control}(\Omega)=
 \langle P,\mathcal X,\mathcal Q,\mathcal N,
          \Gamma,\mathcal E,\mathcal D,w,n\rangle.
\end{equation}
Let $\mathcal R_h\subseteq\dom(\mathcal S_\Omega)$ be the finite set of strong
references temporarily retained by the embedding host, such as arguments of a
platform primitive that is atomic with respect to collection.  A closed
Newspeak execution has $\mathcal R_h=\varnothing$.  First erase the weak
components of a payload while preserving its class and ordinary fields:
\begin{equation}\label{eq:gc-strong-payload}
\begin{aligned}
 \mathsf{strongPayload}(\kw{weakArray}\langle C,\bar y\rangle)
   &=\kw{weakArray}\langle C,\bar{\nil}\rangle,\\
 \mathsf{strongPayload}(\kw{weakMap}\langle C,E_w\rangle)
   &=\kw{weakMap}\langle C,\varnothing\rangle,\\
 \mathsf{strongPayload}(r)&=r
   \quad\text{for every other runtime record }r.
\end{aligned}
\end{equation}
The direct roots, unconditional strong successors, and map-value successors
enabled by a candidate live set $L$ are
\begin{equation}\label{eq:gc-roots-and-successors}
\begin{aligned}
 \mathsf{roots}_{\Omega}(\mathcal R_h)
   &=\bigl(\mathcal R_h\cup
       \mathsf{refs}(\mathsf{control}(\Omega))\bigr)
       \cap\dom(\mathcal S_\Omega),\\
 \mathsf{strongSucc}_{\Omega}(x)
   &=\mathsf{refs}(\mathsf{strongPayload}(
        \mathsf{payload}_{\Omega}(x)))
       \cap\dom(\mathcal S_\Omega),\\
 \mathsf{enabledMapValues}_{\Omega}(x,L)
   &=\begin{cases}
      \displaystyle
      \bigcup_{\substack{k\in\dom(E_w)\\ k\in L}}
        \mathsf{refs}(E_w(k))\cap\dom(\mathcal S_\Omega)
       & \mathsf{payload}_{\Omega}(x)=
           \kw{weakMap}\langle C,E_w\rangle,\\[3mm]
      \varnothing & \text{otherwise.}
     \end{cases}
\end{aligned}
\end{equation}
The structural extractor $\mathsf{refs}$ is Equation~
\eqref{eq:structural-reference-extractor}.  Static platform objects and atoms
are roots through $P$; running activations and their terms are roots through
$\mathcal X$; queued and in-flight arguments, reply promises, and settlement
values are roots through $\mathcal Q$ and $\mathcal N$.  Event identities and
actor identities that do not occur in $\dom(\mathcal S_\Omega)$ are discarded
by the intersection.

Strong reachability is the least fixed point
\begin{equation}\label{eq:gc-reachability}
\begin{aligned}
 \mathsf{reach}^{0}_{\Omega}(\mathcal R_h)
   &=\mathsf{roots}_{\Omega}(\mathcal R_h),\\
 \mathsf{reach}^{i+1}_{\Omega}(\mathcal R_h)
   &=\mathsf{reach}^{i}_{\Omega}(\mathcal R_h)
     \cup\bigcup_{x\in\mathsf{reach}^{i}_{\Omega}(\mathcal R_h)}
       \bigl(\mathsf{strongSucc}_{\Omega}(x)
        \cup\mathsf{enabledMapValues}_{\Omega}
          (x,\mathsf{reach}^{i}_{\Omega}(\mathcal R_h))\bigr),\\
 \mathsf{live}_{\Omega}(\mathcal R_h)
   &=\bigcup_{i\in\mathbb N}\mathsf{reach}^{i}_{\Omega}(\mathcal R_h),\\
 \mathsf{garbage}_{\Omega}(\mathcal R_h)
   &=\dom(\mathcal S_\Omega)
       \setminus\mathsf{live}_{\Omega}(\mathcal R_h).
\end{aligned}
\end{equation}
All stores are finite, so this sequence reaches a fixed point.  Cycles not
reachable from a root are garbage together; reference counting is neither
assumed nor implied.  Weak-array elements never enter the fixed point merely
because their array is live.  A weak-map value enters only after both its map
and its key have entered; consequently a value that points back to its otherwise
unreachable key does not activate its own entry.

\subsection{Admissible collection and identity history}

A collector need not reclaim all eligible records.  It may reclaim a nonempty
set $G$, provided no retained record would contain a dangling strong reference:
\begin{equation}\label{eq:admissible-garbage-set}
\begin{split}
 \mathsf{admissibleGC}_{\Omega,\mathcal R_h}(G)
 \quad\Longleftrightarrow\quad{}
 &\varnothing\ne G\subseteq
     \mathsf{garbage}_{\Omega}(\mathcal R_h)\ \land\\
 &\forall x\in\dom(\mathcal S_\Omega)\setminus G.\;
      \mathsf{strongSucc}_{\Omega}(x)\cap G=\varnothing.
\end{split}
\end{equation}
The second conjunct matters because $\kw{allInstancesOfQuery}$ can expose a
retained but previously unreachable object.  It rules out retaining such an
object while collecting another object to which one of its ordinary slots
still points.  Taking all of $\mathsf{garbage}_{\Omega}(\mathcal R_h)$ always
satisfies this condition; a generational or incremental collector may choose a
smaller closed subset.

Weak edges are intentionally absent from the closure premise.  If a selected
record is named by a retained weak array or by an inactive weak-map entry, the
collector clears that weak observation rather than rejecting the collection.
An enabled weak-map value cannot belong to $G$ because Equation~
\eqref{eq:gc-reachability} already makes it live.

Define weak-edge clearing on a retained payload by
\begin{equation}\label{eq:gc-clear-weak-payload}
\begin{aligned}
 \mathsf{clearWeak}_{G}
  (\kw{weakArray}\langle C,\langle y_1,\ldots,y_m\rangle\rangle)
  &=\kw{weakArray}\langle C,\langle y'_1,\ldots,y'_m\rangle\rangle,\\
 y'_i&=\begin{cases}
       \nil&\mathsf{refs}(y_i)\cap G\ne\varnothing,\\
       y_i&\mathsf{refs}(y_i)\cap G=\varnothing,
      \end{cases}\\
 \mathsf{clearWeak}_{G}(\kw{weakMap}\langle C,E_w\rangle)
  &=\kw{weakMap}\langle C,E'_w\rangle,\\
 E'_w&=E_w|_{\{k\mid k\notin G\ \land\
                    \mathsf{refs}(E_w(k))\cap G=\varnothing\}},\\
 \mathsf{clearWeak}_{G}(r)&=r
   \quad\text{for every other runtime record }r.
\end{aligned}
\end{equation}
The value-side test on $E'_w$ permits partial collection to remove an inactive
entry's value while leaving its still-unreachable key for a later collection;
no retained weak record then contains a dangling identity.

For $K=\dom(\mathcal S_\Omega)\setminus G$, let
\begin{equation}\label{eq:gc-pruned-actor-heap}
 \mathsf{prunedActorHeap}(H_A,G,K)=
 \{x\mapsto\mathsf{clearWeak}_{G}(H_A(x))
      \mid x\in\dom(H_A)\cap K\}.
\end{equation}
Store pruning is the total operation
\begin{equation}\label{eq:gc-prune-store}
\begin{split}
 \mathsf{prune}(\Omega,G)=\Omega[&
   V\mapsto V|_K,\quad
   \mathcal H(A)\mapsto
      \mathsf{prunedActorHeap}(\mathcal H(A),G,K)
      \text{ for every }A\in\dom(\mathcal H),\\
  &\Pi\mapsto\Pi|_K,\quad\Phi\mapsto\Phi|_K].
\end{split}
\end{equation}
Map restriction is by key.  Every other component, including causal and event
history, is unchanged.

Deleting a record must not make its identity fresh again.  Rather than adding
an allocation ledger to every existing configuration and rule, the optional
layer wraps the VM state:
\begin{equation}\label{eq:gc-wrapper-state}
 \Psi=\langle\Omega,\mathcal I\rangle,
 \qquad
 \dom(\mathcal S_\Omega)\subseteq\mathcal I\subseteq\mathit{Ref},
\end{equation}
where $\mathcal I$ is the finite set of reference identities allocated so far
in this execution.  Let the combined existing VM transition be the disjoint
union
\begin{equation}\label{eq:gc-underlying-transition}
\begin{aligned}
 \Omega\mathrel{\Longrightarrow^{\kw{actorStep}}}\Omega'
  &\Longleftrightarrow \Omega\longrightarrow_{\mathsf{act}}\Omega',\\
 \Omega\mathrel{\Longrightarrow^{\lambda_r}}\Omega'
  &\Longleftrightarrow \Omega\xrightarrow{\lambda_r}\Omega',
 \quad
 \lambda_r::=\kw{commitReflection}(A,\Delta,n)
       \mid\kw{rejectReflection}(A,\Delta,\xi_r).
\end{aligned}
\end{equation}
For one such transition define
\begin{equation}\label{eq:gc-new-identities}
 \mathsf{newRefs}(\Omega,\Omega')=
   \dom(\mathcal S_{\Omega'})\setminus\dom(\mathcal S_{\Omega}).
\end{equation}
The ordinary transition lifts only when all newly allocated identities have
never been used:
\[
\dfrac{
 \Omega\mathrel{\Longrightarrow^{\lambda}}\Omega'
 \qquad
 \mathsf{newRefs}(\Omega,\Omega')\cap\mathcal I=\varnothing
}{
 \mathcal R_h\vdash\langle\Omega,\mathcal I\rangle
   \mathrel{\Longrightarrow^{\lambda}_{\mathsf{gc}}}
 \langle\Omega',
   \mathcal I\cup\mathsf{newRefs}(\Omega,\Omega')\rangle
}
\quad(\rulelabel{GC-Lift})
\]
This gives every existing ``fresh'' premise its intended global meaning while
leaving the underlying rules unchanged.  An initial wrapped configuration has
$\mathcal I=\dom(\mathcal S_{\Omega_0})$.

Collection is an independently enabled VM-safepoint transition:
\[
\dfrac{
 \mathsf{admissibleGC}_{\Omega,\mathcal R_h}(G)
}{
 \mathcal R_h\vdash\langle\Omega,\mathcal I\rangle
   \longrightarrow_{\mathsf{gc}}
 \langle\mathsf{prune}(\Omega,G),\mathcal I\rangle
}
\quad(\rulelabel{GC-Collect})
\]
There is deliberately no fairness requirement for
\rulelabel{GC-Collect}: an eligible object may remain uncollected indefinitely.
The implementation chooses both when to collect and which admissible closed
garbage region to reclaim.

\subsection{Observable consequences}

Let $\Omega'=\mathsf{prune}(\Omega,G)$.  Equations~
\eqref{eq:gc-reachability}--\eqref{eq:gc-prune-store} immediately give the
safety obligations
\begin{equation}\label{eq:gc-safety}
\begin{aligned}
 \mathsf{live}_{\Omega}(\mathcal R_h)
   &\subseteq\dom(\mathcal S_{\Omega'}),\\
 \forall x\in\dom(\mathcal S_{\Omega'}).\quad
 \mathsf{refs}(\mathsf{payload}_{\Omega'}(x))
   \cap\mathit{Ref}
   &\subseteq\dom(\mathcal S_{\Omega'}).
\end{aligned}
\end{equation}
The first property says that no strongly reachable record is reclaimed.  The
second says that every retained encoded reference still resolves: selected
weak targets have already been replaced by $\nil$ or removed from their map.
Together
with Equation~\eqref{eq:gc-wrapper-state}, they state memory safety and
permanent identity freshness independently of any particular tracing
algorithm.

For exact-instance enumeration, collection has the intentionally observable
effect
\begin{equation}\label{eq:gc-all-instances-effect}
 \mathsf{allInstancesOf}_{\Omega'}(C)=
 \mathsf{filter}_{x\notin G}
   \bigl(\mathsf{allInstancesOf}_{\Omega}(C)\bigr).
\end{equation}
For mixin-application enumeration, Equation~\eqref{eq:mixin-applications}
similarly gives
\begin{equation}\label{eq:gc-mixin-applications-effect}
 \mathsf{mixinApplications}_{\Omega'}(M)=
 \mathsf{filter}_{C\notin G}
   \bigl(\mathsf{mixinApplications}_{\Omega}(M)\bigr).
\end{equation}
Let $H_A=\mathcal H(A)$ and
$H'_A=\mathsf{prunedActorHeap}(H_A,G,K)$.  For a retained weak array,
Equation~\eqref{eq:gc-clear-weak-payload} gives the observable read
\begin{equation}\label{eq:gc-weak-array-effect}
\begin{aligned}
 y&=\mathsf{weakArrayRead}(H_A,a_w,i),\qquad a_w\in K,\\
 \mathsf{weakArrayRead}(H'_A,a_w,i)
  &=\begin{cases}
     \nil & \mathsf{refs}(y)\cap G\ne\varnothing,\\
     y & \text{otherwise.}
    \end{cases}
\end{aligned}
\end{equation}
For a retained weak map, an entry disappears when collection removes its key
or any reference in its value:
\begin{equation}\label{eq:gc-weak-map-effect}
\begin{aligned}
 v&=\mathsf{weakMapLookup}(H_A,m_w,k),\qquad m_w\in K,\\
 \mathsf{weakMapLookup}(H'_A,m_w,k)
  &=\begin{cases}
     \nil & k\in G\ \lor\ \mathsf{refs}(v)\cap G\ne\varnothing,\\
     v & \text{otherwise.}
    \end{cases}
\end{aligned}
\end{equation}
These changes occur atomically with \rulelabel{GC-Collect}; no program step can
observe a cleared weak cell while the corresponding target record is still
being removed.  A strongly reachable key keeps its weak-map value live through
Equation~\eqref{eq:gc-reachability}, while neither the map nor its value keeps
an otherwise unreachable key live.
Consequently an unreachable instance or class application can appear in a
snapshot before a GC step and be absent after it.  Mirror inspection itself is
atomic with respect to \rulelabel{GC-Collect}; once a returned sequence is
installed in the caller's activation or heap, its elements are ordinary strong
references and are roots of the next collection through $\mathcal X$ or a heap
record.  Nothing can be resurrected after its record has been removed.

If programs have no GC-observing primitive, removing unreachable records is
expected to be a stuttering refinement of the uncollected semantics.  In this
semantics the population queries deliberately invalidate that premise, so the
timing of collection is part of the permitted observable nondeterminism.
Finalizers, if later added to the Newspeak platform, must extend the collection
rule here rather than modify the base evaluator.
| line restores the monotonically growing heaps of
Sections~3--11.  Including it replaces that implementation abstraction with an
explicit, nondeterministic collector.  The collector introduces no Newspeak
syntax and never runs inside an actor step or reflective transaction.

The construction follows the explicit-heap rewrite method of Morrisett,
Felleisen, and Harper, \emph{Abstract Models of Memory Management} (1995), and
the treatment of observable collection in Donnelly, Hallett, and Kfoury,
\emph{Formal Semantics of Weak References} (2006).  The distinctive Newspeak
observations are Equation~\eqref{eq:all-instances-of}, whose implicit instance
registry is weak, and the platform classes \kw{WeakArray} and \kw{WeakMap}.
Collection can therefore change both later population snapshots and later
reads of weak containers.

\subsection{The collectable store}

Garbage collection applies to ordinary heap records and to the routing records
of promise and far-reference handles.  For
$\Omega=\langle P,V,\mathcal H,\mathcal X,\mathcal Q,\mathcal N,
\Pi,\Phi,\Gamma,\mathcal E,\mathcal D,w,n\rangle$, define the disjoint reference
store
\begin{equation}\label{eq:gc-reference-store}
 \mathsf{refStore}(\Omega)
   =\mathsf{vmStore}(\Omega)\uplus\Pi\uplus\Phi.
\end{equation}
The union is disjoint because its three domains are respectively subsets of
$\mathit{ObjRef}$, $\mathit{PromiseId}$, and $\mathit{FarRefId}$.  Put
$\mathcal S_\Omega=\mathsf{refStore}(\Omega)$ and expose a uniform payload
projection:
\begin{equation}\label{eq:gc-reference-payload}
 \mathsf{payload}_{\Omega}(x)=
 \begin{cases}
  \mathsf{vmStore}(\Omega)(x) & x\in\mathit{ObjRef},\\
  \Pi(x) & x\in\mathit{PromiseId},\\
  \Phi(x) & x\in\mathit{FarRefId}.
 \end{cases}
 \qquad x\in\dom(\mathcal S_\Omega).
\end{equation}
Thus a live promise strongly retains its settlement value and registered
waiters, and a live far-reference handle strongly retains its near target.
The tables $\Pi$ and $\Phi$ do not retain their own keys merely by containing
them; otherwise every handle and every remote target would live forever.

\subsection{WeakArray and WeakMap}
\label{sec:weak-containers}

\kw{WeakArray} and \kw{WeakMap} are mutable platform objects and therefore
belong to exactly one actor heap.  Their concrete VM layouts are abstracted by
two additional object-record forms:
\begin{equation}\label{eq:weak-container-records}
\begin{aligned}
 H_A(a_w)&=\kw{weakArray}\langle C,\langle y_1,\ldots,y_m\rangle\rangle,
   &a_w&\in\mathit{WeakArrayId},\\
 H_A(m_w)&=\kw{weakMap}\langle C,E_w\rangle,
   &m_w&\in\mathit{WeakMapId},
\end{aligned}
\end{equation}
where $\mathit{WeakArrayId}$ and $\mathit{WeakMapId}$ are disjoint subsets of
$\mathit{ObjRef}$, every $y_i\in\mathit{Ref}$, and
$E_w:\mathit{ObjRef}\rightharpoonup\mathit{Ref}$ is finite.  $C$ and any
ordinary fixed fields of the implementation are strong; the sequence elements
and map keys are weak.  A map value is conditionally strong exactly while its
key is strongly reachable.  This is the usual ephemeron interpretation of a
weak-key map: a value-to-key cycle cannot keep the key alive.

Every stored weak identity initially resolves; weakness controls retention, not
whether an arbitrary dangling identity may be written:
\begin{equation}\label{eq:weak-container-validity}
\begin{aligned}
 \mathsf{weakArrayStoreOK}(\Omega)
 &\Longleftrightarrow
   \forall A,a_w,C,\bar y.\
     \mathcal H(A)(a_w)=\kw{weakArray}\langle C,\bar y\rangle\\[-1mm]
 &\hspace{35mm}\Longrightarrow
     \mathsf{refs}(\bar y)\subseteq\dom(\mathcal S_\Omega),\\
 \mathsf{weakMapStoreOK}(\Omega)
 &\Longleftrightarrow
   \forall A,m_w,C,E_w.\
     \mathcal H(A)(m_w)=\kw{weakMap}\langle C,E_w\rangle\\[-1mm]
 &\hspace{35mm}\Longrightarrow
     \dom(E_w)\cup\mathsf{refs}(\mathsf{range}(E_w))
       \subseteq\dom(\mathcal S_\Omega),\\
 \mathsf{weakContainersOK}(\Omega)
 &\Longleftrightarrow
   \mathsf{weakArrayStoreOK}(\Omega)\land\mathsf{weakMapStoreOK}(\Omega).
\end{aligned}
\end{equation}

The platform method contract is given by distinct partial store functions
rather than overloaded notation:
\begin{equation}\label{eq:weak-array-operations}
\begin{aligned}
 \mathsf{weakArrayRead}(H_A,a_w,i)&=y_i,\\
 \mathsf{writeWeakArrayCell}(H_A,a_w,i,v)
   &=H_A[a_w\mapsto
       \kw{weakArray}\langle C,\bar y[i\mapsto v]\rangle],\\
 &\hspace{35mm}1\le i\le |\bar y|.
\end{aligned}
\end{equation}
\begin{equation}\label{eq:weak-map-operations}
\begin{aligned}
 \mathsf{weakMapLookup}(H_A,m_w,k)&=
   \begin{cases}E_w(k)&k\in\dom(E_w),\\ \nil&k\notin\dom(E_w),\end{cases}\\
 \mathsf{writeWeakMapEntry}(H_A,m_w,k,v)
   &=H_A[m_w\mapsto\kw{weakMap}\langle C,E_w[k\mapsto v]\rangle].
\end{aligned}
\end{equation}
Here $H_A(a_w)$ and $H_A(m_w)$ have the respective forms of
Equation~\eqref{eq:weak-container-records}.  Allocation creates an all-$\nil$
array of the requested nonnegative size or a map with $E_w=\varnothing$, using
the ordinary global-freshness obligation of Equation~\eqref{eq:gc-wrapper-state}.
The two writes are defined only when their new key and value references resolve
in $\mathcal S_\Omega$, and they preserve Equation~
\eqref{eq:weak-container-validity}.
Out-of-bounds array indices are undefined and are reified by the platform
exception boundary.  Storing $\nil$ is permitted.  Since lookup also returns
$\nil$ for an absent key, the public \kw{WeakMap} protocol deliberately does
not distinguish absence from a stored $\nil$ value.

The state outside the collectable store is
\begin{equation}\label{eq:gc-control-state}
 \mathsf{control}(\Omega)=
 \langle P,\mathcal X,\mathcal Q,\mathcal N,
          \Gamma,\mathcal E,\mathcal D,w,n\rangle.
\end{equation}
Let $\mathcal R_h\subseteq\dom(\mathcal S_\Omega)$ be the finite set of strong
references temporarily retained by the embedding host, such as arguments of a
platform primitive that is atomic with respect to collection.  A closed
Newspeak execution has $\mathcal R_h=\varnothing$.  First erase the weak
components of a payload while preserving its class and ordinary fields:
\begin{equation}\label{eq:gc-strong-payload}
\begin{aligned}
 \mathsf{strongPayload}(\kw{weakArray}\langle C,\bar y\rangle)
   &=\kw{weakArray}\langle C,\bar{\nil}\rangle,\\
 \mathsf{strongPayload}(\kw{weakMap}\langle C,E_w\rangle)
   &=\kw{weakMap}\langle C,\varnothing\rangle,\\
 \mathsf{strongPayload}(r)&=r
   \quad\text{for every other runtime record }r.
\end{aligned}
\end{equation}
The direct roots, unconditional strong successors, and map-value successors
enabled by a candidate live set $L$ are
\begin{equation}\label{eq:gc-roots-and-successors}
\begin{aligned}
 \mathsf{roots}_{\Omega}(\mathcal R_h)
   &=\bigl(\mathcal R_h\cup
       \mathsf{refs}(\mathsf{control}(\Omega))\bigr)
       \cap\dom(\mathcal S_\Omega),\\
 \mathsf{strongSucc}_{\Omega}(x)
   &=\mathsf{refs}(\mathsf{strongPayload}(
        \mathsf{payload}_{\Omega}(x)))
       \cap\dom(\mathcal S_\Omega),\\
 \mathsf{enabledMapValues}_{\Omega}(x,L)
   &=\begin{cases}
      \displaystyle
      \bigcup_{\substack{k\in\dom(E_w)\\ k\in L}}
        \mathsf{refs}(E_w(k))\cap\dom(\mathcal S_\Omega)
       & \mathsf{payload}_{\Omega}(x)=
           \kw{weakMap}\langle C,E_w\rangle,\\[3mm]
      \varnothing & \text{otherwise.}
     \end{cases}
\end{aligned}
\end{equation}
The structural extractor $\mathsf{refs}$ is Equation~
\eqref{eq:structural-reference-extractor}.  Static platform objects and atoms
are roots through $P$; running activations and their terms are roots through
$\mathcal X$; queued and in-flight arguments, reply promises, and settlement
values are roots through $\mathcal Q$ and $\mathcal N$.  Event identities and
actor identities that do not occur in $\dom(\mathcal S_\Omega)$ are discarded
by the intersection.

Strong reachability is the least fixed point
\begin{equation}\label{eq:gc-reachability}
\begin{aligned}
 \mathsf{reach}^{0}_{\Omega}(\mathcal R_h)
   &=\mathsf{roots}_{\Omega}(\mathcal R_h),\\
 \mathsf{reach}^{i+1}_{\Omega}(\mathcal R_h)
   &=\mathsf{reach}^{i}_{\Omega}(\mathcal R_h)
     \cup\bigcup_{x\in\mathsf{reach}^{i}_{\Omega}(\mathcal R_h)}
       \bigl(\mathsf{strongSucc}_{\Omega}(x)
        \cup\mathsf{enabledMapValues}_{\Omega}
          (x,\mathsf{reach}^{i}_{\Omega}(\mathcal R_h))\bigr),\\
 \mathsf{live}_{\Omega}(\mathcal R_h)
   &=\bigcup_{i\in\mathbb N}\mathsf{reach}^{i}_{\Omega}(\mathcal R_h),\\
 \mathsf{garbage}_{\Omega}(\mathcal R_h)
   &=\dom(\mathcal S_\Omega)
       \setminus\mathsf{live}_{\Omega}(\mathcal R_h).
\end{aligned}
\end{equation}
All stores are finite, so this sequence reaches a fixed point.  Cycles not
reachable from a root are garbage together; reference counting is neither
assumed nor implied.  Weak-array elements never enter the fixed point merely
because their array is live.  A weak-map value enters only after both its map
and its key have entered; consequently a value that points back to its otherwise
unreachable key does not activate its own entry.

\subsection{Admissible collection and identity history}

A collector need not reclaim all eligible records.  It may reclaim a nonempty
set $G$, provided no retained record would contain a dangling strong reference:
\begin{equation}\label{eq:admissible-garbage-set}
\begin{split}
 \mathsf{admissibleGC}_{\Omega,\mathcal R_h}(G)
 \quad\Longleftrightarrow\quad{}
 &\varnothing\ne G\subseteq
     \mathsf{garbage}_{\Omega}(\mathcal R_h)\ \land\\
 &\forall x\in\dom(\mathcal S_\Omega)\setminus G.\;
      \mathsf{strongSucc}_{\Omega}(x)\cap G=\varnothing.
\end{split}
\end{equation}
The second conjunct matters because $\kw{allInstancesOfQuery}$ can expose a
retained but previously unreachable object.  It rules out retaining such an
object while collecting another object to which one of its ordinary slots
still points.  Taking all of $\mathsf{garbage}_{\Omega}(\mathcal R_h)$ always
satisfies this condition; a generational or incremental collector may choose a
smaller closed subset.

Weak edges are intentionally absent from the closure premise.  If a selected
record is named by a retained weak array or by an inactive weak-map entry, the
collector clears that weak observation rather than rejecting the collection.
An enabled weak-map value cannot belong to $G$ because Equation~
\eqref{eq:gc-reachability} already makes it live.

Define weak-edge clearing on a retained payload by
\begin{equation}\label{eq:gc-clear-weak-payload}
\begin{aligned}
 \mathsf{clearWeak}_{G}
  (\kw{weakArray}\langle C,\langle y_1,\ldots,y_m\rangle\rangle)
  &=\kw{weakArray}\langle C,\langle y'_1,\ldots,y'_m\rangle\rangle,\\
 y'_i&=\begin{cases}
       \nil&\mathsf{refs}(y_i)\cap G\ne\varnothing,\\
       y_i&\mathsf{refs}(y_i)\cap G=\varnothing,
      \end{cases}\\
 \mathsf{clearWeak}_{G}(\kw{weakMap}\langle C,E_w\rangle)
  &=\kw{weakMap}\langle C,E'_w\rangle,\\
 E'_w&=E_w|_{\{k\mid k\notin G\ \land\
                    \mathsf{refs}(E_w(k))\cap G=\varnothing\}},\\
 \mathsf{clearWeak}_{G}(r)&=r
   \quad\text{for every other runtime record }r.
\end{aligned}
\end{equation}
The value-side test on $E'_w$ permits partial collection to remove an inactive
entry's value while leaving its still-unreachable key for a later collection;
no retained weak record then contains a dangling identity.

For $K=\dom(\mathcal S_\Omega)\setminus G$, let
\begin{equation}\label{eq:gc-pruned-actor-heap}
 \mathsf{prunedActorHeap}(H_A,G,K)=
 \{x\mapsto\mathsf{clearWeak}_{G}(H_A(x))
      \mid x\in\dom(H_A)\cap K\}.
\end{equation}
Store pruning is the total operation
\begin{equation}\label{eq:gc-prune-store}
\begin{split}
 \mathsf{prune}(\Omega,G)=\Omega[&
   V\mapsto V|_K,\quad
   \mathcal H(A)\mapsto
      \mathsf{prunedActorHeap}(\mathcal H(A),G,K)
      \text{ for every }A\in\dom(\mathcal H),\\
  &\Pi\mapsto\Pi|_K,\quad\Phi\mapsto\Phi|_K].
\end{split}
\end{equation}
Map restriction is by key.  Every other component, including causal and event
history, is unchanged.

Deleting a record must not make its identity fresh again.  Rather than adding
an allocation ledger to every existing configuration and rule, the optional
layer wraps the VM state:
\begin{equation}\label{eq:gc-wrapper-state}
 \Psi=\langle\Omega,\mathcal I\rangle,
 \qquad
 \dom(\mathcal S_\Omega)\subseteq\mathcal I\subseteq\mathit{Ref},
\end{equation}
where $\mathcal I$ is the finite set of reference identities allocated so far
in this execution.  Let the combined existing VM transition be the disjoint
union
\begin{equation}\label{eq:gc-underlying-transition}
\begin{aligned}
 \Omega\mathrel{\Longrightarrow^{\kw{actorStep}}}\Omega'
  &\Longleftrightarrow \Omega\longrightarrow_{\mathsf{act}}\Omega',\\
 \Omega\mathrel{\Longrightarrow^{\lambda_r}}\Omega'
  &\Longleftrightarrow \Omega\xrightarrow{\lambda_r}\Omega',
 \quad
 \lambda_r::=\kw{commitReflection}(A,\Delta,n)
       \mid\kw{rejectReflection}(A,\Delta,\xi_r).
\end{aligned}
\end{equation}
For one such transition define
\begin{equation}\label{eq:gc-new-identities}
 \mathsf{newRefs}(\Omega,\Omega')=
   \dom(\mathcal S_{\Omega'})\setminus\dom(\mathcal S_{\Omega}).
\end{equation}
The ordinary transition lifts only when all newly allocated identities have
never been used:
\[
\dfrac{
 \Omega\mathrel{\Longrightarrow^{\lambda}}\Omega'
 \qquad
 \mathsf{newRefs}(\Omega,\Omega')\cap\mathcal I=\varnothing
}{
 \mathcal R_h\vdash\langle\Omega,\mathcal I\rangle
   \mathrel{\Longrightarrow^{\lambda}_{\mathsf{gc}}}
 \langle\Omega',
   \mathcal I\cup\mathsf{newRefs}(\Omega,\Omega')\rangle
}
\quad(\rulelabel{GC-Lift})
\]
This gives every existing ``fresh'' premise its intended global meaning while
leaving the underlying rules unchanged.  An initial wrapped configuration has
$\mathcal I=\dom(\mathcal S_{\Omega_0})$.

Collection is an independently enabled VM-safepoint transition:
\[
\dfrac{
 \mathsf{admissibleGC}_{\Omega,\mathcal R_h}(G)
}{
 \mathcal R_h\vdash\langle\Omega,\mathcal I\rangle
   \longrightarrow_{\mathsf{gc}}
 \langle\mathsf{prune}(\Omega,G),\mathcal I\rangle
}
\quad(\rulelabel{GC-Collect})
\]
There is deliberately no fairness requirement for
\rulelabel{GC-Collect}: an eligible object may remain uncollected indefinitely.
The implementation chooses both when to collect and which admissible closed
garbage region to reclaim.

\subsection{Observable consequences}

Let $\Omega'=\mathsf{prune}(\Omega,G)$.  Equations~
\eqref{eq:gc-reachability}--\eqref{eq:gc-prune-store} immediately give the
safety obligations
\begin{equation}\label{eq:gc-safety}
\begin{aligned}
 \mathsf{live}_{\Omega}(\mathcal R_h)
   &\subseteq\dom(\mathcal S_{\Omega'}),\\
 \forall x\in\dom(\mathcal S_{\Omega'}).\quad
 \mathsf{refs}(\mathsf{payload}_{\Omega'}(x))
   \cap\mathit{Ref}
   &\subseteq\dom(\mathcal S_{\Omega'}).
\end{aligned}
\end{equation}
The first property says that no strongly reachable record is reclaimed.  The
second says that every retained encoded reference still resolves: selected
weak targets have already been replaced by $\nil$ or removed from their map.
Together
with Equation~\eqref{eq:gc-wrapper-state}, they state memory safety and
permanent identity freshness independently of any particular tracing
algorithm.

For exact-instance enumeration, collection has the intentionally observable
effect
\begin{equation}\label{eq:gc-all-instances-effect}
 \mathsf{allInstancesOf}_{\Omega'}(C)=
 \mathsf{filter}_{x\notin G}
   \bigl(\mathsf{allInstancesOf}_{\Omega}(C)\bigr).
\end{equation}
For mixin-application enumeration, Equation~\eqref{eq:mixin-applications}
similarly gives
\begin{equation}\label{eq:gc-mixin-applications-effect}
 \mathsf{mixinApplications}_{\Omega'}(M)=
 \mathsf{filter}_{C\notin G}
   \bigl(\mathsf{mixinApplications}_{\Omega}(M)\bigr).
\end{equation}
Let $H_A=\mathcal H(A)$ and
$H'_A=\mathsf{prunedActorHeap}(H_A,G,K)$.  For a retained weak array,
Equation~\eqref{eq:gc-clear-weak-payload} gives the observable read
\begin{equation}\label{eq:gc-weak-array-effect}
\begin{aligned}
 y&=\mathsf{weakArrayRead}(H_A,a_w,i),\qquad a_w\in K,\\
 \mathsf{weakArrayRead}(H'_A,a_w,i)
  &=\begin{cases}
     \nil & \mathsf{refs}(y)\cap G\ne\varnothing,\\
     y & \text{otherwise.}
    \end{cases}
\end{aligned}
\end{equation}
For a retained weak map, an entry disappears when collection removes its key
or any reference in its value:
\begin{equation}\label{eq:gc-weak-map-effect}
\begin{aligned}
 v&=\mathsf{weakMapLookup}(H_A,m_w,k),\qquad m_w\in K,\\
 \mathsf{weakMapLookup}(H'_A,m_w,k)
  &=\begin{cases}
     \nil & k\in G\ \lor\ \mathsf{refs}(v)\cap G\ne\varnothing,\\
     v & \text{otherwise.}
    \end{cases}
\end{aligned}
\end{equation}
These changes occur atomically with \rulelabel{GC-Collect}; no program step can
observe a cleared weak cell while the corresponding target record is still
being removed.  A strongly reachable key keeps its weak-map value live through
Equation~\eqref{eq:gc-reachability}, while neither the map nor its value keeps
an otherwise unreachable key live.
Consequently an unreachable instance or class application can appear in a
snapshot before a GC step and be absent after it.  Mirror inspection itself is
atomic with respect to \rulelabel{GC-Collect}; once a returned sequence is
installed in the caller's activation or heap, its elements are ordinary strong
references and are roots of the next collection through $\mathcal X$ or a heap
record.  Nothing can be resurrected after its record has been removed.

If programs have no GC-observing primitive, removing unreachable records is
expected to be a stuttering refinement of the uncollected semantics.  In this
semantics the population queries deliberately invalidate that premise, so the
timing of collection is part of the permitted observable nondeterminism.
Finalizers, if later added to the Newspeak platform, must extend the collection
rule here rather than modify the base evaluator.
| line restores the monotonically growing heaps of
Sections~3--11.  Including it replaces that implementation abstraction with an
explicit, nondeterministic collector.  The collector introduces no Newspeak
syntax and never runs inside an actor step or reflective transaction.

The construction follows the explicit-heap rewrite method of Morrisett,
Felleisen, and Harper, \emph{Abstract Models of Memory Management} (1995), and
the treatment of observable collection in Donnelly, Hallett, and Kfoury,
\emph{Formal Semantics of Weak References} (2006).  The distinctive Newspeak
observations are Equation~\eqref{eq:all-instances-of}, whose implicit instance
registry is weak, and the platform classes \kw{WeakArray} and \kw{WeakMap}.
Collection can therefore change both later population snapshots and later
reads of weak containers.

\subsection{The collectable store}

Garbage collection applies to ordinary heap records and to the routing records
of promise and far-reference handles.  For
$\Omega=\langle P,V,\mathcal H,\mathcal X,\mathcal Q,\mathcal N,
\Pi,\Phi,\Gamma,\mathcal E,\mathcal D,w,n\rangle$, define the disjoint reference
store
\begin{equation}\label{eq:gc-reference-store}
 \mathsf{refStore}(\Omega)
   =\mathsf{vmStore}(\Omega)\uplus\Pi\uplus\Phi.
\end{equation}
The union is disjoint because its three domains are respectively subsets of
$\mathit{ObjRef}$, $\mathit{PromiseId}$, and $\mathit{FarRefId}$.  Put
$\mathcal S_\Omega=\mathsf{refStore}(\Omega)$ and expose a uniform payload
projection:
\begin{equation}\label{eq:gc-reference-payload}
 \mathsf{payload}_{\Omega}(x)=
 \begin{cases}
  \mathsf{vmStore}(\Omega)(x) & x\in\mathit{ObjRef},\\
  \Pi(x) & x\in\mathit{PromiseId},\\
  \Phi(x) & x\in\mathit{FarRefId}.
 \end{cases}
 \qquad x\in\dom(\mathcal S_\Omega).
\end{equation}
Thus a live promise strongly retains its settlement value and registered
waiters, and a live far-reference handle strongly retains its near target.
The tables $\Pi$ and $\Phi$ do not retain their own keys merely by containing
them; otherwise every handle and every remote target would live forever.

\subsection{WeakArray and WeakMap}
\label{sec:weak-containers}

\kw{WeakArray} and \kw{WeakMap} are mutable platform objects and therefore
belong to exactly one actor heap.  Their concrete VM layouts are abstracted by
two additional object-record forms:
\begin{equation}\label{eq:weak-container-records}
\begin{aligned}
 H_A(a_w)&=\kw{weakArray}\langle C,\langle y_1,\ldots,y_m\rangle\rangle,
   &a_w&\in\mathit{WeakArrayId},\\
 H_A(m_w)&=\kw{weakMap}\langle C,E_w\rangle,
   &m_w&\in\mathit{WeakMapId},
\end{aligned}
\end{equation}
where $\mathit{WeakArrayId}$ and $\mathit{WeakMapId}$ are disjoint subsets of
$\mathit{ObjRef}$, every $y_i\in\mathit{Ref}$, and
$E_w:\mathit{ObjRef}\rightharpoonup\mathit{Ref}$ is finite.  $C$ and any
ordinary fixed fields of the implementation are strong; the sequence elements
and map keys are weak.  A map value is conditionally strong exactly while its
key is strongly reachable.  This is the usual ephemeron interpretation of a
weak-key map: a value-to-key cycle cannot keep the key alive.

Every stored weak identity initially resolves; weakness controls retention, not
whether an arbitrary dangling identity may be written:
\begin{equation}\label{eq:weak-container-validity}
\begin{aligned}
 \mathsf{weakArrayStoreOK}(\Omega)
 &\Longleftrightarrow
   \forall A,a_w,C,\bar y.\
     \mathcal H(A)(a_w)=\kw{weakArray}\langle C,\bar y\rangle\\[-1mm]
 &\hspace{35mm}\Longrightarrow
     \mathsf{refs}(\bar y)\subseteq\dom(\mathcal S_\Omega),\\
 \mathsf{weakMapStoreOK}(\Omega)
 &\Longleftrightarrow
   \forall A,m_w,C,E_w.\
     \mathcal H(A)(m_w)=\kw{weakMap}\langle C,E_w\rangle\\[-1mm]
 &\hspace{35mm}\Longrightarrow
     \dom(E_w)\cup\mathsf{refs}(\mathsf{range}(E_w))
       \subseteq\dom(\mathcal S_\Omega),\\
 \mathsf{weakContainersOK}(\Omega)
 &\Longleftrightarrow
   \mathsf{weakArrayStoreOK}(\Omega)\land\mathsf{weakMapStoreOK}(\Omega).
\end{aligned}
\end{equation}

The platform method contract is given by distinct partial store functions
rather than overloaded notation:
\begin{equation}\label{eq:weak-array-operations}
\begin{aligned}
 \mathsf{weakArrayRead}(H_A,a_w,i)&=y_i,\\
 \mathsf{writeWeakArrayCell}(H_A,a_w,i,v)
   &=H_A[a_w\mapsto
       \kw{weakArray}\langle C,\bar y[i\mapsto v]\rangle],\\
 &\hspace{35mm}1\le i\le |\bar y|.
\end{aligned}
\end{equation}
\begin{equation}\label{eq:weak-map-operations}
\begin{aligned}
 \mathsf{weakMapLookup}(H_A,m_w,k)&=
   \begin{cases}E_w(k)&k\in\dom(E_w),\\ \nil&k\notin\dom(E_w),\end{cases}\\
 \mathsf{writeWeakMapEntry}(H_A,m_w,k,v)
   &=H_A[m_w\mapsto\kw{weakMap}\langle C,E_w[k\mapsto v]\rangle].
\end{aligned}
\end{equation}
Here $H_A(a_w)$ and $H_A(m_w)$ have the respective forms of
Equation~\eqref{eq:weak-container-records}.  Allocation creates an all-$\nil$
array of the requested nonnegative size or a map with $E_w=\varnothing$, using
the ordinary global-freshness obligation of Equation~\eqref{eq:gc-wrapper-state}.
The two writes are defined only when their new key and value references resolve
in $\mathcal S_\Omega$, and they preserve Equation~
\eqref{eq:weak-container-validity}.
Out-of-bounds array indices are undefined and are reified by the platform
exception boundary.  Storing $\nil$ is permitted.  Since lookup also returns
$\nil$ for an absent key, the public \kw{WeakMap} protocol deliberately does
not distinguish absence from a stored $\nil$ value.

The state outside the collectable store is
\begin{equation}\label{eq:gc-control-state}
 \mathsf{control}(\Omega)=
 \langle P,\mathcal X,\mathcal Q,\mathcal N,
          \Gamma,\mathcal E,\mathcal D,w,n\rangle.
\end{equation}
Let $\mathcal R_h\subseteq\dom(\mathcal S_\Omega)$ be the finite set of strong
references temporarily retained by the embedding host, such as arguments of a
platform primitive that is atomic with respect to collection.  A closed
Newspeak execution has $\mathcal R_h=\varnothing$.  First erase the weak
components of a payload while preserving its class and ordinary fields:
\begin{equation}\label{eq:gc-strong-payload}
\begin{aligned}
 \mathsf{strongPayload}(\kw{weakArray}\langle C,\bar y\rangle)
   &=\kw{weakArray}\langle C,\bar{\nil}\rangle,\\
 \mathsf{strongPayload}(\kw{weakMap}\langle C,E_w\rangle)
   &=\kw{weakMap}\langle C,\varnothing\rangle,\\
 \mathsf{strongPayload}(r)&=r
   \quad\text{for every other runtime record }r.
\end{aligned}
\end{equation}
The direct roots, unconditional strong successors, and map-value successors
enabled by a candidate live set $L$ are
\begin{equation}\label{eq:gc-roots-and-successors}
\begin{aligned}
 \mathsf{roots}_{\Omega}(\mathcal R_h)
   &=\bigl(\mathcal R_h\cup
       \mathsf{refs}(\mathsf{control}(\Omega))\bigr)
       \cap\dom(\mathcal S_\Omega),\\
 \mathsf{strongSucc}_{\Omega}(x)
   &=\mathsf{refs}(\mathsf{strongPayload}(
        \mathsf{payload}_{\Omega}(x)))
       \cap\dom(\mathcal S_\Omega),\\
 \mathsf{enabledMapValues}_{\Omega}(x,L)
   &=\begin{cases}
      \displaystyle
      \bigcup_{\substack{k\in\dom(E_w)\\ k\in L}}
        \mathsf{refs}(E_w(k))\cap\dom(\mathcal S_\Omega)
       & \mathsf{payload}_{\Omega}(x)=
           \kw{weakMap}\langle C,E_w\rangle,\\[3mm]
      \varnothing & \text{otherwise.}
     \end{cases}
\end{aligned}
\end{equation}
The structural extractor $\mathsf{refs}$ is Equation~
\eqref{eq:structural-reference-extractor}.  Static platform objects and atoms
are roots through $P$; running activations and their terms are roots through
$\mathcal X$; queued and in-flight arguments, reply promises, and settlement
values are roots through $\mathcal Q$ and $\mathcal N$.  Event identities and
actor identities that do not occur in $\dom(\mathcal S_\Omega)$ are discarded
by the intersection.

Strong reachability is the least fixed point
\begin{equation}\label{eq:gc-reachability}
\begin{aligned}
 \mathsf{reach}^{0}_{\Omega}(\mathcal R_h)
   &=\mathsf{roots}_{\Omega}(\mathcal R_h),\\
 \mathsf{reach}^{i+1}_{\Omega}(\mathcal R_h)
   &=\mathsf{reach}^{i}_{\Omega}(\mathcal R_h)
     \cup\bigcup_{x\in\mathsf{reach}^{i}_{\Omega}(\mathcal R_h)}
       \bigl(\mathsf{strongSucc}_{\Omega}(x)
        \cup\mathsf{enabledMapValues}_{\Omega}
          (x,\mathsf{reach}^{i}_{\Omega}(\mathcal R_h))\bigr),\\
 \mathsf{live}_{\Omega}(\mathcal R_h)
   &=\bigcup_{i\in\mathbb N}\mathsf{reach}^{i}_{\Omega}(\mathcal R_h),\\
 \mathsf{garbage}_{\Omega}(\mathcal R_h)
   &=\dom(\mathcal S_\Omega)
       \setminus\mathsf{live}_{\Omega}(\mathcal R_h).
\end{aligned}
\end{equation}
All stores are finite, so this sequence reaches a fixed point.  Cycles not
reachable from a root are garbage together; reference counting is neither
assumed nor implied.  Weak-array elements never enter the fixed point merely
because their array is live.  A weak-map value enters only after both its map
and its key have entered; consequently a value that points back to its otherwise
unreachable key does not activate its own entry.

\subsection{Admissible collection and identity history}

A collector need not reclaim all eligible records.  It may reclaim a nonempty
set $G$, provided no retained record would contain a dangling strong reference:
\begin{equation}\label{eq:admissible-garbage-set}
\begin{split}
 \mathsf{admissibleGC}_{\Omega,\mathcal R_h}(G)
 \quad\Longleftrightarrow\quad{}
 &\varnothing\ne G\subseteq
     \mathsf{garbage}_{\Omega}(\mathcal R_h)\ \land\\
 &\forall x\in\dom(\mathcal S_\Omega)\setminus G.\;
      \mathsf{strongSucc}_{\Omega}(x)\cap G=\varnothing.
\end{split}
\end{equation}
The second conjunct matters because $\kw{allInstancesOfQuery}$ can expose a
retained but previously unreachable object.  It rules out retaining such an
object while collecting another object to which one of its ordinary slots
still points.  Taking all of $\mathsf{garbage}_{\Omega}(\mathcal R_h)$ always
satisfies this condition; a generational or incremental collector may choose a
smaller closed subset.

Weak edges are intentionally absent from the closure premise.  If a selected
record is named by a retained weak array or by an inactive weak-map entry, the
collector clears that weak observation rather than rejecting the collection.
An enabled weak-map value cannot belong to $G$ because Equation~
\eqref{eq:gc-reachability} already makes it live.

Define weak-edge clearing on a retained payload by
\begin{equation}\label{eq:gc-clear-weak-payload}
\begin{aligned}
 \mathsf{clearWeak}_{G}
  (\kw{weakArray}\langle C,\langle y_1,\ldots,y_m\rangle\rangle)
  &=\kw{weakArray}\langle C,\langle y'_1,\ldots,y'_m\rangle\rangle,\\
 y'_i&=\begin{cases}
       \nil&\mathsf{refs}(y_i)\cap G\ne\varnothing,\\
       y_i&\mathsf{refs}(y_i)\cap G=\varnothing,
      \end{cases}\\
 \mathsf{clearWeak}_{G}(\kw{weakMap}\langle C,E_w\rangle)
  &=\kw{weakMap}\langle C,E'_w\rangle,\\
 E'_w&=E_w|_{\{k\mid k\notin G\ \land\
                    \mathsf{refs}(E_w(k))\cap G=\varnothing\}},\\
 \mathsf{clearWeak}_{G}(r)&=r
   \quad\text{for every other runtime record }r.
\end{aligned}
\end{equation}
The value-side test on $E'_w$ permits partial collection to remove an inactive
entry's value while leaving its still-unreachable key for a later collection;
no retained weak record then contains a dangling identity.

For $K=\dom(\mathcal S_\Omega)\setminus G$, let
\begin{equation}\label{eq:gc-pruned-actor-heap}
 \mathsf{prunedActorHeap}(H_A,G,K)=
 \{x\mapsto\mathsf{clearWeak}_{G}(H_A(x))
      \mid x\in\dom(H_A)\cap K\}.
\end{equation}
Store pruning is the total operation
\begin{equation}\label{eq:gc-prune-store}
\begin{split}
 \mathsf{prune}(\Omega,G)=\Omega[&
   V\mapsto V|_K,\quad
   \mathcal H(A)\mapsto
      \mathsf{prunedActorHeap}(\mathcal H(A),G,K)
      \text{ for every }A\in\dom(\mathcal H),\\
  &\Pi\mapsto\Pi|_K,\quad\Phi\mapsto\Phi|_K].
\end{split}
\end{equation}
Map restriction is by key.  Every other component, including causal and event
history, is unchanged.

Deleting a record must not make its identity fresh again.  Rather than adding
an allocation ledger to every existing configuration and rule, the optional
layer wraps the VM state:
\begin{equation}\label{eq:gc-wrapper-state}
 \Psi=\langle\Omega,\mathcal I\rangle,
 \qquad
 \dom(\mathcal S_\Omega)\subseteq\mathcal I\subseteq\mathit{Ref},
\end{equation}
where $\mathcal I$ is the finite set of reference identities allocated so far
in this execution.  Let the combined existing VM transition be the disjoint
union
\begin{equation}\label{eq:gc-underlying-transition}
\begin{aligned}
 \Omega\mathrel{\Longrightarrow^{\kw{actorStep}}}\Omega'
  &\Longleftrightarrow \Omega\longrightarrow_{\mathsf{act}}\Omega',\\
 \Omega\mathrel{\Longrightarrow^{\lambda_r}}\Omega'
  &\Longleftrightarrow \Omega\xrightarrow{\lambda_r}\Omega',
 \quad
 \lambda_r::=\kw{commitReflection}(A,\Delta,n)
       \mid\kw{rejectReflection}(A,\Delta,\xi_r).
\end{aligned}
\end{equation}
For one such transition define
\begin{equation}\label{eq:gc-new-identities}
 \mathsf{newRefs}(\Omega,\Omega')=
   \dom(\mathcal S_{\Omega'})\setminus\dom(\mathcal S_{\Omega}).
\end{equation}
The ordinary transition lifts only when all newly allocated identities have
never been used:
\[
\dfrac{
 \Omega\mathrel{\Longrightarrow^{\lambda}}\Omega'
 \qquad
 \mathsf{newRefs}(\Omega,\Omega')\cap\mathcal I=\varnothing
}{
 \mathcal R_h\vdash\langle\Omega,\mathcal I\rangle
   \mathrel{\Longrightarrow^{\lambda}_{\mathsf{gc}}}
 \langle\Omega',
   \mathcal I\cup\mathsf{newRefs}(\Omega,\Omega')\rangle
}
\quad(\rulelabel{GC-Lift})
\]
This gives every existing ``fresh'' premise its intended global meaning while
leaving the underlying rules unchanged.  An initial wrapped configuration has
$\mathcal I=\dom(\mathcal S_{\Omega_0})$.

Collection is an independently enabled VM-safepoint transition:
\[
\dfrac{
 \mathsf{admissibleGC}_{\Omega,\mathcal R_h}(G)
}{
 \mathcal R_h\vdash\langle\Omega,\mathcal I\rangle
   \longrightarrow_{\mathsf{gc}}
 \langle\mathsf{prune}(\Omega,G),\mathcal I\rangle
}
\quad(\rulelabel{GC-Collect})
\]
There is deliberately no fairness requirement for
\rulelabel{GC-Collect}: an eligible object may remain uncollected indefinitely.
The implementation chooses both when to collect and which admissible closed
garbage region to reclaim.

\subsection{Observable consequences}

Let $\Omega'=\mathsf{prune}(\Omega,G)$.  Equations~
\eqref{eq:gc-reachability}--\eqref{eq:gc-prune-store} immediately give the
safety obligations
\begin{equation}\label{eq:gc-safety}
\begin{aligned}
 \mathsf{live}_{\Omega}(\mathcal R_h)
   &\subseteq\dom(\mathcal S_{\Omega'}),\\
 \forall x\in\dom(\mathcal S_{\Omega'}).\quad
 \mathsf{refs}(\mathsf{payload}_{\Omega'}(x))
   \cap\mathit{Ref}
   &\subseteq\dom(\mathcal S_{\Omega'}).
\end{aligned}
\end{equation}
The first property says that no strongly reachable record is reclaimed.  The
second says that every retained encoded reference still resolves: selected
weak targets have already been replaced by $\nil$ or removed from their map.
Together
with Equation~\eqref{eq:gc-wrapper-state}, they state memory safety and
permanent identity freshness independently of any particular tracing
algorithm.

For exact-instance enumeration, collection has the intentionally observable
effect
\begin{equation}\label{eq:gc-all-instances-effect}
 \mathsf{allInstancesOf}_{\Omega'}(C)=
 \mathsf{filter}_{x\notin G}
   \bigl(\mathsf{allInstancesOf}_{\Omega}(C)\bigr).
\end{equation}
For mixin-application enumeration, Equation~\eqref{eq:mixin-applications}
similarly gives
\begin{equation}\label{eq:gc-mixin-applications-effect}
 \mathsf{mixinApplications}_{\Omega'}(M)=
 \mathsf{filter}_{C\notin G}
   \bigl(\mathsf{mixinApplications}_{\Omega}(M)\bigr).
\end{equation}
Let $H_A=\mathcal H(A)$ and
$H'_A=\mathsf{prunedActorHeap}(H_A,G,K)$.  For a retained weak array,
Equation~\eqref{eq:gc-clear-weak-payload} gives the observable read
\begin{equation}\label{eq:gc-weak-array-effect}
\begin{aligned}
 y&=\mathsf{weakArrayRead}(H_A,a_w,i),\qquad a_w\in K,\\
 \mathsf{weakArrayRead}(H'_A,a_w,i)
  &=\begin{cases}
     \nil & \mathsf{refs}(y)\cap G\ne\varnothing,\\
     y & \text{otherwise.}
    \end{cases}
\end{aligned}
\end{equation}
For a retained weak map, an entry disappears when collection removes its key
or any reference in its value:
\begin{equation}\label{eq:gc-weak-map-effect}
\begin{aligned}
 v&=\mathsf{weakMapLookup}(H_A,m_w,k),\qquad m_w\in K,\\
 \mathsf{weakMapLookup}(H'_A,m_w,k)
  &=\begin{cases}
     \nil & k\in G\ \lor\ \mathsf{refs}(v)\cap G\ne\varnothing,\\
     v & \text{otherwise.}
    \end{cases}
\end{aligned}
\end{equation}
These changes occur atomically with \rulelabel{GC-Collect}; no program step can
observe a cleared weak cell while the corresponding target record is still
being removed.  A strongly reachable key keeps its weak-map value live through
Equation~\eqref{eq:gc-reachability}, while neither the map nor its value keeps
an otherwise unreachable key live.
Consequently an unreachable instance or class application can appear in a
snapshot before a GC step and be absent after it.  Mirror inspection itself is
atomic with respect to \rulelabel{GC-Collect}; once a returned sequence is
installed in the caller's activation or heap, its elements are ordinary strong
references and are roots of the next collection through $\mathcal X$ or a heap
record.  Nothing can be resurrected after its record has been removed.

If programs have no GC-observing primitive, removing unreachable records is
expected to be a stuttering refinement of the uncollected semantics.  In this
semantics the population queries deliberately invalidate that premise, so the
timing of collection is part of the permitted observable nondeterminism.
Finalizers, if later added to the Newspeak platform, must extend the collection
rule here rather than modify the base evaluator.
| line restores the monotonically growing heaps of
Sections~3--11.  Including it replaces that implementation abstraction with an
explicit, nondeterministic collector.  The collector introduces no Newspeak
syntax and never runs inside an actor step or reflective transaction.

The construction follows the explicit-heap rewrite method of Morrisett,
Felleisen, and Harper, \emph{Abstract Models of Memory Management} (1995), and
the treatment of observable collection in Donnelly, Hallett, and Kfoury,
\emph{Formal Semantics of Weak References} (2006).  The distinctive Newspeak
observations are Equation~\eqref{eq:all-instances-of}, whose implicit instance
registry is weak, and the platform classes \kw{WeakArray} and \kw{WeakMap}.
Collection can therefore change both later population snapshots and later
reads of weak containers.

\subsection{The collectable store}

Garbage collection applies to ordinary heap records and to the routing records
of promise and far-reference handles.  For
$\Omega=\langle P,V,\mathcal H,\mathcal X,\mathcal Q,\mathcal N,
\Pi,\Phi,\Gamma,\mathcal E,\mathcal D,w,n\rangle$, define the disjoint reference
store
\begin{equation}\label{eq:gc-reference-store}
 \mathsf{refStore}(\Omega)
   =\mathsf{vmStore}(\Omega)\uplus\Pi\uplus\Phi.
\end{equation}
The union is disjoint because its three domains are respectively subsets of
$\mathit{ObjRef}$, $\mathit{PromiseId}$, and $\mathit{FarRefId}$.  Put
$\mathcal S_\Omega=\mathsf{refStore}(\Omega)$ and expose a uniform payload
projection:
\begin{equation}\label{eq:gc-reference-payload}
 \mathsf{payload}_{\Omega}(x)=
 \begin{cases}
  \mathsf{vmStore}(\Omega)(x) & x\in\mathit{ObjRef},\\
  \Pi(x) & x\in\mathit{PromiseId},\\
  \Phi(x) & x\in\mathit{FarRefId}.
 \end{cases}
 \qquad x\in\dom(\mathcal S_\Omega).
\end{equation}
Thus a live promise strongly retains its settlement value and registered
waiters, and a live far-reference handle strongly retains its near target.
The tables $\Pi$ and $\Phi$ do not retain their own keys merely by containing
them; otherwise every handle and every remote target would live forever.

\subsection{WeakArray and WeakMap}
\label{sec:weak-containers}

\kw{WeakArray} and \kw{WeakMap} are mutable platform objects and therefore
belong to exactly one actor heap.  Their concrete VM layouts are abstracted by
two additional object-record forms:
\begin{equation}\label{eq:weak-container-records}
\begin{aligned}
 H_A(a_w)&=\kw{weakArray}\langle C,\langle y_1,\ldots,y_m\rangle\rangle,
   &a_w&\in\mathit{WeakArrayId},\\
 H_A(m_w)&=\kw{weakMap}\langle C,E_w\rangle,
   &m_w&\in\mathit{WeakMapId},
\end{aligned}
\end{equation}
where $\mathit{WeakArrayId}$ and $\mathit{WeakMapId}$ are disjoint subsets of
$\mathit{ObjRef}$, every $y_i\in\mathit{Ref}$, and
$E_w:\mathit{ObjRef}\rightharpoonup\mathit{Ref}$ is finite.  $C$ and any
ordinary fixed fields of the implementation are strong; the sequence elements
and map keys are weak.  A map value is conditionally strong exactly while its
key is strongly reachable.  This is the usual ephemeron interpretation of a
weak-key map: a value-to-key cycle cannot keep the key alive.

Every stored weak identity initially resolves; weakness controls retention, not
whether an arbitrary dangling identity may be written:
\begin{equation}\label{eq:weak-container-validity}
\begin{aligned}
 \mathsf{weakArrayStoreOK}(\Omega)
 &\Longleftrightarrow
   \forall A,a_w,C,\bar y.\
     \mathcal H(A)(a_w)=\kw{weakArray}\langle C,\bar y\rangle\\[-1mm]
 &\hspace{35mm}\Longrightarrow
     \mathsf{refs}(\bar y)\subseteq\dom(\mathcal S_\Omega),\\
 \mathsf{weakMapStoreOK}(\Omega)
 &\Longleftrightarrow
   \forall A,m_w,C,E_w.\
     \mathcal H(A)(m_w)=\kw{weakMap}\langle C,E_w\rangle\\[-1mm]
 &\hspace{35mm}\Longrightarrow
     \dom(E_w)\cup\mathsf{refs}(\mathsf{range}(E_w))
       \subseteq\dom(\mathcal S_\Omega),\\
 \mathsf{weakContainersOK}(\Omega)
 &\Longleftrightarrow
   \mathsf{weakArrayStoreOK}(\Omega)\land\mathsf{weakMapStoreOK}(\Omega).
\end{aligned}
\end{equation}

The platform method contract is given by distinct partial store functions
rather than overloaded notation:
\begin{equation}\label{eq:weak-array-operations}
\begin{aligned}
 \mathsf{weakArrayRead}(H_A,a_w,i)&=y_i,\\
 \mathsf{writeWeakArrayCell}(H_A,a_w,i,v)
   &=H_A[a_w\mapsto
       \kw{weakArray}\langle C,\bar y[i\mapsto v]\rangle],\\
 &\hspace{35mm}1\le i\le |\bar y|.
\end{aligned}
\end{equation}
\begin{equation}\label{eq:weak-map-operations}
\begin{aligned}
 \mathsf{weakMapLookup}(H_A,m_w,k)&=
   \begin{cases}E_w(k)&k\in\dom(E_w),\\ \nil&k\notin\dom(E_w),\end{cases}\\
 \mathsf{writeWeakMapEntry}(H_A,m_w,k,v)
   &=H_A[m_w\mapsto\kw{weakMap}\langle C,E_w[k\mapsto v]\rangle].
\end{aligned}
\end{equation}
Here $H_A(a_w)$ and $H_A(m_w)$ have the respective forms of
Equation~\eqref{eq:weak-container-records}.  Allocation creates an all-$\nil$
array of the requested nonnegative size or a map with $E_w=\varnothing$, using
the ordinary global-freshness obligation of Equation~\eqref{eq:gc-wrapper-state}.
The two writes are defined only when their new key and value references resolve
in $\mathcal S_\Omega$, and they preserve Equation~
\eqref{eq:weak-container-validity}.
Out-of-bounds array indices are undefined and are reified by the platform
exception boundary.  Storing $\nil$ is permitted.  Since lookup also returns
$\nil$ for an absent key, the public \kw{WeakMap} protocol deliberately does
not distinguish absence from a stored $\nil$ value.

The state outside the collectable store is
\begin{equation}\label{eq:gc-control-state}
 \mathsf{control}(\Omega)=
 \langle P,\mathcal X,\mathcal Q,\mathcal N,
          \Gamma,\mathcal E,\mathcal D,w,n\rangle.
\end{equation}
Let $\mathcal R_h\subseteq\dom(\mathcal S_\Omega)$ be the finite set of strong
references temporarily retained by the embedding host, such as arguments of a
platform primitive that is atomic with respect to collection.  A closed
Newspeak execution has $\mathcal R_h=\varnothing$.  First erase the weak
components of a payload while preserving its class and ordinary fields:
\begin{equation}\label{eq:gc-strong-payload}
\begin{aligned}
 \mathsf{strongPayload}(\kw{weakArray}\langle C,\bar y\rangle)
   &=\kw{weakArray}\langle C,\bar{\nil}\rangle,\\
 \mathsf{strongPayload}(\kw{weakMap}\langle C,E_w\rangle)
   &=\kw{weakMap}\langle C,\varnothing\rangle,\\
 \mathsf{strongPayload}(r)&=r
   \quad\text{for every other runtime record }r.
\end{aligned}
\end{equation}
The direct roots, unconditional strong successors, and map-value successors
enabled by a candidate live set $L$ are
\begin{equation}\label{eq:gc-roots-and-successors}
\begin{aligned}
 \mathsf{roots}_{\Omega}(\mathcal R_h)
   &=\bigl(\mathcal R_h\cup
       \mathsf{refs}(\mathsf{control}(\Omega))\bigr)
       \cap\dom(\mathcal S_\Omega),\\
 \mathsf{strongSucc}_{\Omega}(x)
   &=\mathsf{refs}(\mathsf{strongPayload}(
        \mathsf{payload}_{\Omega}(x)))
       \cap\dom(\mathcal S_\Omega),\\
 \mathsf{enabledMapValues}_{\Omega}(x,L)
   &=\begin{cases}
      \displaystyle
      \bigcup_{\substack{k\in\dom(E_w)\\ k\in L}}
        \mathsf{refs}(E_w(k))\cap\dom(\mathcal S_\Omega)
       & \mathsf{payload}_{\Omega}(x)=
           \kw{weakMap}\langle C,E_w\rangle,\\[3mm]
      \varnothing & \text{otherwise.}
     \end{cases}
\end{aligned}
\end{equation}
The structural extractor $\mathsf{refs}$ is Equation~
\eqref{eq:structural-reference-extractor}.  Static platform objects and atoms
are roots through $P$; running activations and their terms are roots through
$\mathcal X$; queued and in-flight arguments, reply promises, and settlement
values are roots through $\mathcal Q$ and $\mathcal N$.  Event identities and
actor identities that do not occur in $\dom(\mathcal S_\Omega)$ are discarded
by the intersection.

Strong reachability is the least fixed point
\begin{equation}\label{eq:gc-reachability}
\begin{aligned}
 \mathsf{reach}^{0}_{\Omega}(\mathcal R_h)
   &=\mathsf{roots}_{\Omega}(\mathcal R_h),\\
 \mathsf{reach}^{i+1}_{\Omega}(\mathcal R_h)
   &=\mathsf{reach}^{i}_{\Omega}(\mathcal R_h)
     \cup\bigcup_{x\in\mathsf{reach}^{i}_{\Omega}(\mathcal R_h)}
       \bigl(\mathsf{strongSucc}_{\Omega}(x)
        \cup\mathsf{enabledMapValues}_{\Omega}
          (x,\mathsf{reach}^{i}_{\Omega}(\mathcal R_h))\bigr),\\
 \mathsf{live}_{\Omega}(\mathcal R_h)
   &=\bigcup_{i\in\mathbb N}\mathsf{reach}^{i}_{\Omega}(\mathcal R_h),\\
 \mathsf{garbage}_{\Omega}(\mathcal R_h)
   &=\dom(\mathcal S_\Omega)
       \setminus\mathsf{live}_{\Omega}(\mathcal R_h).
\end{aligned}
\end{equation}
All stores are finite, so this sequence reaches a fixed point.  Cycles not
reachable from a root are garbage together; reference counting is neither
assumed nor implied.  Weak-array elements never enter the fixed point merely
because their array is live.  A weak-map value enters only after both its map
and its key have entered; consequently a value that points back to its otherwise
unreachable key does not activate its own entry.

\subsection{Admissible collection and identity history}

A collector need not reclaim all eligible records.  It may reclaim a nonempty
set $G$, provided no retained record would contain a dangling strong reference:
\begin{equation}\label{eq:admissible-garbage-set}
\begin{split}
 \mathsf{admissibleGC}_{\Omega,\mathcal R_h}(G)
 \quad\Longleftrightarrow\quad{}
 &\varnothing\ne G\subseteq
     \mathsf{garbage}_{\Omega}(\mathcal R_h)\ \land\\
 &\forall x\in\dom(\mathcal S_\Omega)\setminus G.\;
      \mathsf{strongSucc}_{\Omega}(x)\cap G=\varnothing.
\end{split}
\end{equation}
The second conjunct matters because $\kw{allInstancesOfQuery}$ can expose a
retained but previously unreachable object.  It rules out retaining such an
object while collecting another object to which one of its ordinary slots
still points.  Taking all of $\mathsf{garbage}_{\Omega}(\mathcal R_h)$ always
satisfies this condition; a generational or incremental collector may choose a
smaller closed subset.

Weak edges are intentionally absent from the closure premise.  If a selected
record is named by a retained weak array or by an inactive weak-map entry, the
collector clears that weak observation rather than rejecting the collection.
An enabled weak-map value cannot belong to $G$ because Equation~
\eqref{eq:gc-reachability} already makes it live.

Define weak-edge clearing on a retained payload by
\begin{equation}\label{eq:gc-clear-weak-payload}
\begin{aligned}
 \mathsf{clearWeak}_{G}
  (\kw{weakArray}\langle C,\langle y_1,\ldots,y_m\rangle\rangle)
  &=\kw{weakArray}\langle C,\langle y'_1,\ldots,y'_m\rangle\rangle,\\
 y'_i&=\begin{cases}
       \nil&\mathsf{refs}(y_i)\cap G\ne\varnothing,\\
       y_i&\mathsf{refs}(y_i)\cap G=\varnothing,
      \end{cases}\\
 \mathsf{clearWeak}_{G}(\kw{weakMap}\langle C,E_w\rangle)
  &=\kw{weakMap}\langle C,E'_w\rangle,\\
 E'_w&=E_w|_{\{k\mid k\notin G\ \land\
                    \mathsf{refs}(E_w(k))\cap G=\varnothing\}},\\
 \mathsf{clearWeak}_{G}(r)&=r
   \quad\text{for every other runtime record }r.
\end{aligned}
\end{equation}
The value-side test on $E'_w$ permits partial collection to remove an inactive
entry's value while leaving its still-unreachable key for a later collection;
no retained weak record then contains a dangling identity.

For $K=\dom(\mathcal S_\Omega)\setminus G$, let
\begin{equation}\label{eq:gc-pruned-actor-heap}
 \mathsf{prunedActorHeap}(H_A,G,K)=
 \{x\mapsto\mathsf{clearWeak}_{G}(H_A(x))
      \mid x\in\dom(H_A)\cap K\}.
\end{equation}
Store pruning is the total operation
\begin{equation}\label{eq:gc-prune-store}
\begin{split}
 \mathsf{prune}(\Omega,G)=\Omega[&
   V\mapsto V|_K,\quad
   \mathcal H(A)\mapsto
      \mathsf{prunedActorHeap}(\mathcal H(A),G,K)
      \text{ for every }A\in\dom(\mathcal H),\\
  &\Pi\mapsto\Pi|_K,\quad\Phi\mapsto\Phi|_K].
\end{split}
\end{equation}
Map restriction is by key.  Every other component, including causal and event
history, is unchanged.

Deleting a record must not make its identity fresh again.  Rather than adding
an allocation ledger to every existing configuration and rule, the optional
layer wraps the VM state:
\begin{equation}\label{eq:gc-wrapper-state}
 \Psi=\langle\Omega,\mathcal I\rangle,
 \qquad
 \dom(\mathcal S_\Omega)\subseteq\mathcal I\subseteq\mathit{Ref},
\end{equation}
where $\mathcal I$ is the finite set of reference identities allocated so far
in this execution.  Let the combined existing VM transition be the disjoint
union
\begin{equation}\label{eq:gc-underlying-transition}
\begin{aligned}
 \Omega\mathrel{\Longrightarrow^{\kw{actorStep}}}\Omega'
  &\Longleftrightarrow \Omega\longrightarrow_{\mathsf{act}}\Omega',\\
 \Omega\mathrel{\Longrightarrow^{\lambda_r}}\Omega'
  &\Longleftrightarrow \Omega\xrightarrow{\lambda_r}\Omega',
 \quad
 \lambda_r::=\kw{commitReflection}(A,\Delta,n)
       \mid\kw{rejectReflection}(A,\Delta,\xi_r).
\end{aligned}
\end{equation}
For one such transition define
\begin{equation}\label{eq:gc-new-identities}
 \mathsf{newRefs}(\Omega,\Omega')=
   \dom(\mathcal S_{\Omega'})\setminus\dom(\mathcal S_{\Omega}).
\end{equation}
The ordinary transition lifts only when all newly allocated identities have
never been used:
\[
\dfrac{
 \Omega\mathrel{\Longrightarrow^{\lambda}}\Omega'
 \qquad
 \mathsf{newRefs}(\Omega,\Omega')\cap\mathcal I=\varnothing
}{
 \mathcal R_h\vdash\langle\Omega,\mathcal I\rangle
   \mathrel{\Longrightarrow^{\lambda}_{\mathsf{gc}}}
 \langle\Omega',
   \mathcal I\cup\mathsf{newRefs}(\Omega,\Omega')\rangle
}
\quad(\rulelabel{GC-Lift})
\]
This gives every existing ``fresh'' premise its intended global meaning while
leaving the underlying rules unchanged.  An initial wrapped configuration has
$\mathcal I=\dom(\mathcal S_{\Omega_0})$.

Collection is an independently enabled VM-safepoint transition:
\[
\dfrac{
 \mathsf{admissibleGC}_{\Omega,\mathcal R_h}(G)
}{
 \mathcal R_h\vdash\langle\Omega,\mathcal I\rangle
   \longrightarrow_{\mathsf{gc}}
 \langle\mathsf{prune}(\Omega,G),\mathcal I\rangle
}
\quad(\rulelabel{GC-Collect})
\]
There is deliberately no fairness requirement for
\rulelabel{GC-Collect}: an eligible object may remain uncollected indefinitely.
The implementation chooses both when to collect and which admissible closed
garbage region to reclaim.

\subsection{Observable consequences}

Let $\Omega'=\mathsf{prune}(\Omega,G)$.  Equations~
\eqref{eq:gc-reachability}--\eqref{eq:gc-prune-store} immediately give the
safety obligations
\begin{equation}\label{eq:gc-safety}
\begin{aligned}
 \mathsf{live}_{\Omega}(\mathcal R_h)
   &\subseteq\dom(\mathcal S_{\Omega'}),\\
 \forall x\in\dom(\mathcal S_{\Omega'}).\quad
 \mathsf{refs}(\mathsf{payload}_{\Omega'}(x))
   \cap\mathit{Ref}
   &\subseteq\dom(\mathcal S_{\Omega'}).
\end{aligned}
\end{equation}
The first property says that no strongly reachable record is reclaimed.  The
second says that every retained encoded reference still resolves: selected
weak targets have already been replaced by $\nil$ or removed from their map.
Together
with Equation~\eqref{eq:gc-wrapper-state}, they state memory safety and
permanent identity freshness independently of any particular tracing
algorithm.

For exact-instance enumeration, collection has the intentionally observable
effect
\begin{equation}\label{eq:gc-all-instances-effect}
 \mathsf{allInstancesOf}_{\Omega'}(C)=
 \mathsf{filter}_{x\notin G}
   \bigl(\mathsf{allInstancesOf}_{\Omega}(C)\bigr).
\end{equation}
For mixin-application enumeration, Equation~\eqref{eq:mixin-applications}
similarly gives
\begin{equation}\label{eq:gc-mixin-applications-effect}
 \mathsf{mixinApplications}_{\Omega'}(M)=
 \mathsf{filter}_{C\notin G}
   \bigl(\mathsf{mixinApplications}_{\Omega}(M)\bigr).
\end{equation}
Let $H_A=\mathcal H(A)$ and
$H'_A=\mathsf{prunedActorHeap}(H_A,G,K)$.  For a retained weak array,
Equation~\eqref{eq:gc-clear-weak-payload} gives the observable read
\begin{equation}\label{eq:gc-weak-array-effect}
\begin{aligned}
 y&=\mathsf{weakArrayRead}(H_A,a_w,i),\qquad a_w\in K,\\
 \mathsf{weakArrayRead}(H'_A,a_w,i)
  &=\begin{cases}
     \nil & \mathsf{refs}(y)\cap G\ne\varnothing,\\
     y & \text{otherwise.}
    \end{cases}
\end{aligned}
\end{equation}
For a retained weak map, an entry disappears when collection removes its key
or any reference in its value:
\begin{equation}\label{eq:gc-weak-map-effect}
\begin{aligned}
 v&=\mathsf{weakMapLookup}(H_A,m_w,k),\qquad m_w\in K,\\
 \mathsf{weakMapLookup}(H'_A,m_w,k)
  &=\begin{cases}
     \nil & k\in G\ \lor\ \mathsf{refs}(v)\cap G\ne\varnothing,\\
     v & \text{otherwise.}
    \end{cases}
\end{aligned}
\end{equation}
These changes occur atomically with \rulelabel{GC-Collect}; no program step can
observe a cleared weak cell while the corresponding target record is still
being removed.  A strongly reachable key keeps its weak-map value live through
Equation~\eqref{eq:gc-reachability}, while neither the map nor its value keeps
an otherwise unreachable key live.
Consequently an unreachable instance or class application can appear in a
snapshot before a GC step and be absent after it.  Mirror inspection itself is
atomic with respect to \rulelabel{GC-Collect}; once a returned sequence is
installed in the caller's activation or heap, its elements are ordinary strong
references and are roots of the next collection through $\mathcal X$ or a heap
record.  Nothing can be resurrected after its record has been removed.

If programs have no GC-observing primitive, removing unreachable records is
expected to be a stuttering refinement of the uncollected semantics.  In this
semantics the population queries deliberately invalidate that premise, so the
timing of collection is part of the permitted observable nondeterminism.
Finalizers, if later added to the Newspeak platform, must extend the collection
rule here rather than modify the base evaluator.
